\documentclass{article}
\usepackage{comment}
\usepackage{fullpage}
\usepackage{subcaption}
%\usepackage{minipage}
\usepackage{graphicx}
%\usepackage{enumerate}
\usepackage{enumitem}
%\usepackage{calc}
\usepackage{xcolor}
\usepackage[sort&compress,numbers]{natbib}
\usepackage[colorlinks=true, linkcolor=blue, citecolor=blue, urlcolor=blue]{hyperref}
\usepackage{xfrac}
\usepackage{rotating}

\usepackage{amsthm}
\usepackage{amsmath}
\usepackage{amsfonts}
\usepackage{amssymb}
\usepackage{euscript}
\usepackage{bm}
\usepackage{pdflscape} % Add in preamble
\usepackage{framed}
\usepackage{graphicx}
\usepackage{tikz}
\usetikzlibrary{arrows.meta, positioning}
\usepackage{mathrsfs}  
\usepackage[affil-it]{authblk}
\usepackage{dsfont}


\makeatletter
\def\namedlabel#1#2{\begingroup
  #2%
  \def\@currentlabel{#2}%
  \phantomsection\label{#1}\endgroup
}
\makeatother
\newcommand{\caselabel}[2]{\hypertarget{#1}{#2}}
\newcommand{\caseref}[1]{\hyperlink{#1}{#1}}



\renewcommand{\S}{\mathbf{S}}
\newcommand{\C}{\mathbf{C}}

\newcommand{\bA}{\mathbf{A}}
\newcommand{\bZ}{\mathbf{Z}}
\newcommand{\bB}{\mathbf{B}}
\newcommand{\bM}{\mathbf{M}}
\newcommand{\bP}{\mathbf{P}}
\newcommand{\bU}{\mathbf{U}}
\newcommand{\bX}{\mathbf{X}}
\newcommand{\bT}{\mathbf{T}}

\newcommand{\bSigma}{\mathbf{\Sigma}}

\newcommand{\bx}{\mathbf{x}}
\newcommand{\bu}{\mathbf{u}}
\newcommand{\bv}{\mathbf{v}}
\renewcommand{\bf}{\mathbf{f}}
\renewcommand{\bm}{\mathbf{m}}


\newcommand{\sC}{\EuScript{C}}
\newcommand{\sH}{\EuScript{H}}
\newcommand{\sN}{\EuScript{N}}
\newcommand{\sR}{\EuScript{R}}
\newcommand{\sS}{\EuScript{S}}

\newcommand{\bbP}{\mathbb{P}}
\newcommand{\bbQ}{\mathbb{Q}}
\newcommand{\bbN}{\mathbb{N}}
\newcommand{\bbR}{\mathbb{R}}
\newcommand{\bbC}{\mathbb{C}}

\newcommand{\bzero}{\mathbf{0}}

\DeclareMathOperator{\Ker}{Ker}
\DeclareMathOperator{\Ran}{Ran}

\newcommand{\sA}{\EuScript{A}}
\newcommand{\sB}{\EuScript{B}}

\newcommand{\sD}{\EuScript{D}}
\newcommand{\sE}{\EuScript{E}}
\newcommand{\sF}{\EuScript{F}}
\newcommand{\sG}{\EuScript{G}}
\newcommand{\sI}{\EuScript{I}}
\newcommand{\sJ}{\EuScript{J}}
\newcommand{\sT}{\EuScript{T}}
\newcommand{\sQ}{\EuScript{Q}}

\newcommand{\sX}{\EuScript{X}}

\newcommand{\balpha}{\bm{\alpha}}



\newcommand{\bK}{\mathbf{K}}
\newcommand{\bJ}{\mathbf{J}}
\newcommand{\bS}{\mathbf{S}}
\newcommand{\bH}{\mathbf{H}}
\newcommand{\bC}{\mathbf{C}}

\newcommand{\cF}{\mathcal{F}}
\newcommand{\cP}{\mathcal{P}}
\newcommand{\cS}{\mathcal{S}}
\newcommand{\cT}{\mathcal{T}}
\newcommand{\cH}{\mathcal{H}}
\newcommand{\cI}{\mathcal{I}}
\newcommand{\cX}{\mathcal{X}}
\newcommand{\cA}{\mathcal{A}}
\newcommand{\cB}{\mathcal{B}}
\newcommand{\cL}{\mathcal{L}}
\newcommand{\cN}{\mathcal{N}}
\newcommand{\cR}{\mathcal{R}}
\newcommand{\cC}{\mathcal{C}}
\newcommand{\cK}{\mathcal{K}}


\newcommand{\bI}{\mathbf{I}}
\newcommand{\m}{\mathbf{m}}

\newcommand{\bbE}{\mathbb{E}}
\renewcommand{\bbP}{\mathbb{P}}
\newcommand{\bbZ}{\mathbb{Z}}


\newcommand{\bbGamma}{{\mathpalette\makebbGamma\relax}}
\newcommand{\makebbGamma}[2]{%
  \raisebox{\depth}{\scalebox{1}[-1]{$\mathsurround=0pt#1\mathbb{L}$}}%
}




%MATH OPERATORS
\DeclareMathOperator*{\argmax}{arg\,max}
\DeclareMathOperator*{\argmin}{arg\,min}
%\DeclareMathOperator{\argmax}{argmax}
%\DeclareMathOperator{\argmin}{argmin}
%\DeclareMathOperator{\limsup}{limsup}

\DeclareMathOperator{\tr}{tr}
\DeclareMathOperator{\dg}{dg}
\DeclareMathOperator{\rank}{rank}
\DeclareMathOperator{\supp}{supp}
\DeclareMathOperator{\Span}{Span}

\DeclareMathOperator{\Ext}{Ext}
\DeclareMathOperator{\glim}{\Gamma-lim}

\DeclareMathOperator{\opi}{\hat{\otimes}_{\pi}}

%STAT OPERATORS
\DeclareMathOperator{\Cov}{Cov}
\newcommand{\CI}{\mathrel{\text{\scalebox{1.07}{$\perp\mkern-10mu\perp$}}}}
%\newcommand{\tlinf}{\mathrel{$|$}-\mathrel{$|$}}
%THEOREM ENVIRONMENTS
\theoremstyle{plain}
\newtheorem{definition}{Definition}[section]
\newtheorem{lemma}[definition]{Lemma}
\newtheorem{theorem}[definition]{Theorem}
\newtheorem{proposition}[definition]{Proposition}
\newtheorem{corollary}[definition]{Corollary}
\newtheorem{remark}[definition]{Remark}
\newtheorem{example}[definition]{Example}
\newtheorem{conjecture}[definition]{Conjecture}
\newtheorem{assumption}{Assumption}
% change representation of \theassumption changes the numbering style
\renewcommand{\theassumption}{\arabic{assumption}}

%Highlighting
\newcommand{\RED}{\textcolor{red}}

% for tables
\usepackage{multirow}
\usepackage{multicol}
\newcommand\Tstrut{\rule{0pt}{2.6ex}}        
\newcommand\Bstrut{\rule[-0.9ex]{0pt}{0pt}}  

% for captions for tables and figures
\usepackage{caption} 
\captionsetup[table]{skip=20pt}

%Commenting
\definecolor{kartik}{rgb}{0.8,0, 0.2}
\newcommand{\kartik}[1]{\textcolor{kartik}{#1}}

\newcommand{\red}[1]{{\color{red} #1}}
\newcommand{\trace}{\operatorname{trace}}
\newcommand{\range}{\operatorname{range}}

\usepackage{xfrac}  

\newcommand{\test}{\mathfrak{T}}
\newcommand{\Id}{\mathbf{I}}


\newcommand{\dettwo}{\operatorname{det}_2}

\newcommand{\mybox}[1]{\begin{framed}\noindent#1\end{framed}}




% table
\setlength{\arrayrulewidth}{0.5mm}
\setlength{\tabcolsep}{7pt}
\renewcommand{\arraystretch}{1.5}

\title{Near Perfect Classification with Kernel Likelihood Ratios}
\author{Leonardo Santoro}
\date{December 2025}

\begin{document}

\maketitle


This research plan stems from a recent line of work \cite{santoro2025from,santoroLRT25} by myself and coauthors, that uncovers novel information-theoretic and structural–geometric foundations of kernel methods, and opens the way to a host of new methods and results. These developments were in turn upheld by my earlier contributions to the theory and methodology for Gaussian measures on Hilbert spaces in the Bures-Wasserstein geometry  \cite{santoro2023large, jaffe2024large} and their dynamical evolution \cite{santoro2023statistical,santoro2024karhunen}, which uphold many analytical and probabilistic tools underlying this framework.


 In our recent preprint \cite{santoro2025from}  we showed that kernel embeddings elicit a peculiar phenomenon of perfect separation between probability distributions, understood in an information-theoretic sense.
 That is, we showed that if $\bbP,\bbQ$ are probability measures on a compact and separable metric space $\cX$, and $k: \cX \times \cX \to \bbR$ is a universal reproducing kernel thereon, then suitable Gaussian measures constructed from the second-order (covariance) kernel embeddings of $\bbP$ and $\bbQ$ are \emph{mutually singular} if and only if the original measures are distinct:
 \begin{equation}
  \label{eq:equivalence}
\bbP \neq \bbQ 
\quad\Longleftrightarrow \quad
\cN(0, \S_\bbP) \perp \cN(0,\S_\bbQ)
\end{equation}
Here, $\perp$ denotes the {mutual singularity} of Gaussian measures, and $\cN(0, \S_\bbP), \cN(0,\S_\bbQ)$ denote the centered Gaussian measures on the RKHS $\cH$ associated to the kernel $k(\cdot,\cdot)$ with covariance operators given by the second-order kernel embeddings \eqref{eq:covemb} of $\bbP$ and $\bbQ$, respectively. In turn, it is readily seen that \eqref{eq:equivalence} also implies as a corollary that the singularity effect is elicited also by \emph{uncentered} Gaussian embeddings: $\bbP \neq \bbQ 
\Longleftrightarrow \cN(\m_\bbP, \S_\bbP) \perp \cN(\m_\bbQ,\S_\bbQ)$.
% , where $\m_\bbP, \m_\bbQ$ are the mean embeddings \eqref{eq:meanemb} of $\bbP, \bbQ$.

Whereas the classical injectivity of the kernel mean embedding guarantees only that distinct measures are mapped to distinct elements of the RKHS \cite{muandet2017kernel}, our result demonstrates a much stronger separation, one that is both geometric and information-theoretic in nature.
This strengthened perspective not only motivates the construction of new statistics designed to capture such separations, but also reveals new geometric insight into the discriminative mechanisms of kernel representations.
 In statistical terms, this entails that testing for equality of two probability measures on a compact and separable metric space is equivalent to testing for singularity between two centered Gaussian measures on an RKHS, defined through the corresponding kernel embeddings of a probability distribution, with the Hilbert space generated by the embedding kernel. 
Crucially, distinguishing between singular Gaussian measures is fundamentally simpler from an information-theoretic perspective than the original nonparametric two-sample problem, particularly in complex or high-dimensional domains, since singular Gaussians are supported on essentially disjoint affine subspaces. Our proof leverages the classical Feldman–Hájek dichotomy \cite{feldman1958equivalence,hajek1958property} and shows that even a small perturbation of a distribution is maximally magnified through its Gaussian embedding. We refer to this effect as the \textit{separation of measure phenomenon}, a blessing of infinite dimensionality made possible by kernel embeddings, which crystallizes in a precise mathematical statement the outstanding empirical effectiveness of the “kernel trick.”


 

Such insight is harnessed in follow-up work \cite{santoroLRT25} to develop a novel, powerful criterion for nonparametric hypothesis testing of distributional equality, i.e.\ a kernel-based nonparametric two-sample test grounded in the  separation of measure phenomenon. The test constructs a likelihood-ratio statistic based on the regularised relative entropy \cite{minh2020infinite,minh2021regularized} between the Gaussian embeddings, which  satisfies a striking “0/$\infty$” law, for vanishing regularisation:
% , it vanishes under the null hypothesis of equal distributions and diverges under the alternative:
$$
  \lim_{\gamma \to 0}  
    D_{\gamma, \mathrm{KL}} \big( \cN(\mathbf{0}, \S_{\bbP}) \:||\:  \cN(\mathbf{0}, \S_{\bbQ}) \big) 
      % =  \lim_{\gamma \to 0}    -  \frac{1}{2}
      %     \log\dettwo \left(\Id + (\S_{\bbQ} + \gamma \Id)^{-\sfrac{1}{2}}(\S_{\bbP}  - \S_{\bbQ})(\S_{\bbQ} + \gamma \Id)^{-\sfrac{1}{2}} \right)
  = \begin{cases}
      0, \: &\textnormal{if} \quad \bbP=\bbQ,\\
      \infty, \: &\textnormal{if} \quad \bbP\neq \bbQ.
  \end{cases}
$$
where $\gamma>0$ is a regularisation parameter ensuring well-posedness.
To make this operational at finite samples, we consider calibration via permutation, and establish both consistency and uniform power guarantees under mild assumptions for suitably decaying regularisation, with increasing sample size. Empirically, the method achieves substantial gains in detection power compared to state-of-the-art approaches, especially in high-dimensional
 and weak-signal regimes, while also unifying and extending prior work \cite{eric2007testing, hagrass2024spectral}.

\medskip


\mybox{
It is my strong belief that such developments unlock a new dimension of kernel methods and point to a great unexplored potential. The separation of measure phenomenon suggests that RKHS embeddings may encode far richer statistical information than previously recognized, with implications not only for two-sample testing but also for broader inference tasks, and multiple research avenues.
  % Much remains to be explored  --  and discovered, and I strongly believe that the separation of measure phenomenon can uncover the way forward in multiple research avenues.
}





\subsubsection{(B) Likelihood-based classification with kernel Gaussian embeddings: theory and scalable methodology}
\subsubsection*{B.1.  Likelihood-based classification with kernel Gaussian embeddings}

The mutual singularity of Gaussian embeddings of distinct probability distributions uncovered in my recent work \cite{santoro2025from} promises a new theoretical framework for probabilistic classification. 
Let $\{\bbP_j\}_{j=1}^C$ denote class-conditional distributions on $\cX$, and consider their Gaussian embeddings $\cN(\m_{\bbP_j}, \S_{\bbP_j})$ in $\cH$. Given that Gaussian measures on $\cH$ are pairwise mutually singular -- i.e., $\cN(\m_{\bbP_j}, \S_{\bbP_j}) \perp \cN(\m_{\bbP_i}, \S_{\bbP_i})$ for $i \neq j$ --  we may simplistically say that the support of each Gaussian approximation of the embedded class distribution occupies a distinct affine subspace of $\cH$. This observation suggests that a likelihood-based classifier defined in $\cH$ can, in principle, perfectly separate the classes. The purpose of this project will be to develop, analyze, and implement such a classifier.
The corresponding classification rule I propose to investigate can be formulated as a maximum-likelihood classifier in the RKHS:
\begin{equation}
\label{eq:multiclass_rule}
\argmin_{1 \le j \le C}
\; \Big\{
  \langle \varphi(\bx) - \m_{\bbP_j},\, (\S_{\bbP_j}+\gamma \Id)^{-1} (\varphi(\bx) - \m_{\bbP_j}) \rangle_{\cH}
  + \log \dettwo(\S_{\bbP_j}+\gamma \Id)
\Big\}, \nonumber
\end{equation}
where $\gamma > 0$ ensures well-definition, whilst controlling overfit: smaller values of $\gamma$ sharpen the effective decision boundary by reducing regularisation, but at the cost of increased numerical instability and potential overfitting. The trade-off between discrimination sharpness and stability therefore requires careful theoretical and empirical study (note the connections with kernel quadratic discriminant analysis, with $\gamma$ acting as a shrinkage parameter). In the idealized regime where the feature-mapped distributions $\varphi_\# \bbP_i$ are exactly Gaussian, the classifier achieves zero asymptotic error, which I have proven as direct consequence of the mutual singularity result.
The central theoretical challenge lies in characterizing the regime in which this Gaussian approximation is valid. Specifically, one must quantify the discrepancy between the true pushforwards $\varphi_\# \bbP_j$ and their Gaussian surrogates, and establish how such deviations impact the induced generalization error. A promising route is to exploit approximate Gaussianity induced by random projections: by \citet{diaconis1984asymptotics} high-dimensional random projections of complex distributions tend to exhibit near-Gaussian behaviour, providing a natural regularisation mechanism against model misspecification. Complementarily, I will explore whether kernel design can yield a similar stabilizing effect. The aim will be to obtain a non-asymptotic theoretical characterization of the generalization error of the proposed classifier, elucidating the operational regime in which RKHS Gaussian embeddings yield reliable probabilistic classification.

 To further boost scalability and stability, I will combine them with {ensemble methods}: train \(B\) independent classifiers \(\{f_b\}_{b=1}^B\) on subsampled batches of size \(m\) and or feature dimension \(M\), and aggregate the corresponding classifiers.
 Such  strategies impose an implicit rank constraint, effectively performing \emph{spectral pruning} that discards low-variance directions, which in turn entails an implicit regularisation that mitigates overfitting. 
\bibliographystyle{plainnat}
\bibliography{bib}

\end{document}

